\section{Introduction}
The project is very exploratory in nature, it was planned to be a strategy for leveraging graph structure to make determinations, it has since evolved into development of a new software pattern that offers greater value than the original plan in terms of broader applicability. This project attempts to determine best way to design the pattern, it remains within same topic and uses graph structure,  REPS introduces a structured approach using a set of rules and relationships between nodes to make determinations. The focus of this thesis is to explore the design space for such a pattern rather than to directly compare it to existing solutions, as optimization and performance benchmarking are not the primary objectives.

This article is structured to follow the chronological progression of the project. After the related work section, the direction of the project shifts noticeably—from evaluating REPS as a competitor to existing predictive systems, to assessing its role as a distinct and interpretable software pattern.
	
Modern \gls{ml} systems often suffer from a lack of interpretability. Developers cannot always explain or verify how a model makes its determinations. This makes it challenging for developers to understand how the system operates internally, verify the results, or ensure the verifiability of the outcomes. This black box nature introduces a barrier to trust and understanding in their software's determinations.

To address these challenges, I present \gls{reps}, a graph-based real-time analysis tool which developers themselves can design to fit their use-case while making use of many state-of-the-art machine learning models. \gls{reps} operates by processing input data, which will be either sensor or event data which will pass through nodes called "event nodes" which generates event data. Each node processes inputs from one or more sources and generates a result that can be used on one or multiple downstream nodes. This system enables real-time decision-making across interconnected components while still allowing the developer to know how it operates.

This project aims to determine whether there are any viable use cases for \gls{reps} and what advantages and disadvantages it may have.
\newpage
\subsection{Problem}\label{sec:RQProblem}
Research Questions:
	\begin{itemize}
		\item RQ0: How should a \gls{reps} be constructed?
		\item RQ1: What advantages and disadvantages does \gls{reps}?
	\end{itemize}
	
\noindent
Research Question mapping:
	\begin{itemize}
		\item  RQ0: The construction principles of \gls{reps} can be derived from the practical experience of designing and implementing the system. This question aims to guide future implementations of similar systems by highlighting architectural decisions, trade-offs, and lessons learned throughout the development process.
		\item RQ1: For \gls{reps} to be a viable candidate in real-world environments, it must offer value over existing solutions. Therefore, this question involves a comparative evaluation in terms of CPU usage, memory consumption, configurability, ease of understanding, and predictive capabilities.
	\end{itemize}